\scanpage{17}
% archive sheet number: 16

Système d'hyperplans (projectifs) réels

% slanted
\margin{Théorie analogue dans le cas affine mais qui revient au cas projectif, \unsure{avec|où} un des $H_i$ \unsure{distingués|distingué}.}

Soit $V$ espace vectoriel réel, $H$ un hyperplan de $V$, on note $D_H$ \inserted{$\subset\mathcal P(V)$} l'ens des trois éléments \deleted{de} $\{H,V'_H,V''_H\}$, où $V'_H$ et $V''_H$ sont les deux demi-espaces ouverts de $V$ déterminés par $H$ ; \deleted{qui forment} alors \deleted{\unsure{$H, V'_+$}} $D_H$ est une partition de $V$ par des ens convexes.

Comme c'est un ensemble à trois éléments, un élément distingué \deleted{$\{$}$H$\deleted{$\}$}, on peut regarder $D_H$ comme un espace vectoriel sur $\mathbb F_3=\{-1,0,1\}$. On a une application canonique

\deleted{\unsure{Soit $V_i$|Soit $V\setminus$}}

\[
\varphi_H:\underset{\substack{\shortparallel\\ V^*/\mathbb R^{*+}}}{\mathbb S}\xrightarrow{\deleted{\varphi_H}}D_H
\]

Soit $(H_i)_{\unsure{i|j|I}\in I}$ une famille finie d'hyperplans de $V$, \uline{distincts}, et considérons l'application

% slanted
\margin{``distincts'' inutile, mais on ne \unsure{gagne} rien essentiellement}

\[
\varphi:S\longrightarrow\prod D_{H_i}=E
\]

définie par les $\varphi_{H_i}$. Les fibres non vides de $\varphi$ forment une partition de $S$ par des \deleted{\unsure{sous}} ens ``convexes'' (i.e. l'image inverse dans $V^*$, \unsure{augm. de} $\{0\}$, est \deleted{convexe} \inserted{un} cône convexe) \uline{\unsure{donc|des|dans}} \uline{connexes} : \inserted{les ``cellules ouvertes''.} Soit

% slanted
\margin{``\unsure{demi convexe|demi-convexe|donc connexes|convexes}'' inutile, sauf lorsque le cône convexe en question est une droite, -- mais si $\bigcap H_i=\{0\}$, ceci n'est pas possible. On supposera \unsure{désormais|donc} $\bigcap H_i=\{0\}$.}

\[
\Sigma\deleted{\unsure{\Sigma}}=\varphi(S)\subset E
\]

l'ens. d'indices de cette partition. C'est une partie de $E$ stable par \inserted{l'} homothétie $-1$ ($-\mathrm{id}_E$), qui correspond au passage \deleted{\unsure{inverse}} d'une fibre

\[
Z^*_{-\alpha}=-Z^*_\alpha\qquad\text{pour }\alpha\in\Sigma
\]

On se propose de déterminer $\overline{Z^*_\alpha}$ pour $i\in I$.
